% 本文件用于替换 main.tex 中问题一的“自适应大邻域搜索”、
% “算法实现与输出文件”和“结果展示与算法分析”相关内容。

\subsection{基于整数修复的自适应大邻域搜索}

广义节约算法只能合并总人数不超过单机容量的两条路线，无法改变已经形成的人员分批，
也难以修正早期LAND机场分配对后续路线组合产生的影响。为此，本文以通过独立校验的
广义节约方案 $S_{\mathrm{init}}^{(1)}$ 为初始解，构造带整数规划修复的自适应大邻域
搜索。该算法保留前述路线评价器和燃油标签搜索，不使用近似距离代替最终航段评价。

\subsubsection{破坏算子}

每次迭代从当前方案中移除若干条完整路线，并将其中人员恢复为待分配需求组。程序实际
实现以下五种破坏算子：
\begin{enumerate}[label=\arabic*.]
  \item \textbf{随机路线移除：}随机选择若干路线，用于保持搜索多样性；
  \item \textbf{相关路线移除：}根据服务设施间距离、出发机场和指定机场人员构造
        相关性，优先联合移除空间上接近的路线；
  \item \textbf{低利用率路线移除：}优先移除座位公里利用率较低的余量路线；
  \item \textbf{单位人员时间最差路线移除：}优先移除飞机时间与载客人数比值较大的
        路线；
  \item \textbf{LAND重指派移除：}优先选择LAND人员较多的路线，使其在修复阶段可以
        重新选择A01、A02或A03。
\end{enumerate}

\subsubsection{候选路线生成与整数规划修复}

设本轮被移除的聚合需求组集合为 $\mathcal G_D$，候选路线集合为
$\mathcal R_D^{\mathrm{cand}}$。每条候选路线 $c$ 同时给定出发机场 $a(c)$、机型
$k(c)$ 和至多两个实际服务设施的访问顺序；技术经停点和加油决策由多标签搜索重新
确定。记 $Q_g^D$ 为需求组 $g$ 的人数，$s_{k(c)}$ 为候选路线容量，
$\vartheta_{cg}$ 为该路线将需求组 $g$ 送达目的设施所需的单人在途时间。

定义整数变量 $\zeta_c$ 为候选路线 $c$ 的使用架次数，$\chi_{gc}$ 为需求组 $g$
分配至候选路线 $c$ 的人数。修复模型满足
\begin{align}
  \sum_{c\in\mathcal R_D^{\mathrm{cand}}}\chi_{gc}
  &=Q_g^D,
  &&g\in\mathcal G_D,\\
  \sum_{g\in\mathcal G_D}\chi_{gc}
  &\le s_{k(c)}\zeta_c,
  &&c\in\mathcal R_D^{\mathrm{cand}}.
\end{align}
若需求组 $g$ 的目的设施不在路线 $c$ 中，或其指定机场与 $a(c)$ 不一致，则固定
$\chi_{gc}=0$；起点为LAND的需求可以分配至任意机场。为避免生成没有实际人员需求的
服务停靠，对于路线 $c$ 中的每个服务设施 $f$，进一步要求
\begin{equation}
  \sum_{g:\,f_g=f}\chi_{gc}\ge \zeta_c.
\end{equation}

修复模型首先最小化邻域内总飞机使用时间
\begin{equation}
  \min \sum_{c\in\mathcal R_D^{\mathrm{cand}}}
  T_c^{\mathrm{air}}\zeta_c.
\end{equation}
在固定一级目标最优值后，再按人员在途时间、架次数和燃油消耗的顺序选择修复方案。
求得聚合分配量后，程序按照人员编号回填到具体架次，并调用内存评价器逐航段检查容量、
油量、首次到达目的设施下机和最多5次海上停靠约束。

\subsubsection{自适应选择与接受准则}

设算子 $\ell$ 在当前统计周期内被调用 $n_{\ell}$ 次，累计奖励为 $\sigma_{\ell}$，
其权重按下式更新：
\begin{equation}
  w_{\ell}^{\mathrm{new}}
  =(1-\rho)w_{\ell}^{\mathrm{old}}
  +\rho\frac{\sigma_{\ell}}{\max\{1,n_{\ell}\}}.
\end{equation}
产生历史最优、改进当前解和被接受候选解分别获得不同奖励。若候选解在词典序意义下
优于当前解则直接接受；否则根据第一个发生变化的目标分量计算相对劣化量，并以
\begin{equation}
  P_{\mathrm{acc}}=\exp\!\left(-\frac{\Delta}{\Theta}\right)
\end{equation}
决定是否接受，温度按 $\Theta\leftarrow\alpha\Theta$ 递减。无论当前解是否暂时
劣化，历史最优解始终单独保存。

\subsubsection{两阶段参数设置}

综合解质量与运行时间，本文采用表~\ref{tab:q1-alns-parameters}所示两阶段配置。第一
阶段快速清除明显的低利用率路线，第二阶段扩大破坏规模以重新组合更多人员；单条候选
路线最多包含两个实际服务设施，但标签算法仍可插入额外技术经停点，且总海上停靠次数
不超过5次。

\begin{table}[htbp]
  \centering
  \caption{问题一两阶段ALNS参数}
  \label{tab:q1-alns-parameters}
  \begin{tabular}{lrrrrr}
    \toprule
    阶段 & 随机种子 & 迭代数 & 每次移除路线数 & 修复时限/s & 总时限/s\\
    \midrule
    快速搜索 & 0 & 60 & 2--3 & 3 & 600\\
    强化搜索 & 1 & 60 & 3--4 & 5 & 900\\
    \bottomrule
  \end{tabular}
\end{table}

\subsection{算法实现与可复现性}

问题一程序中，\texttt{src/solver/baseline.py}实现单设施容量动态规划和LAND分流，
\texttt{src/solver/improve.py}实现广义节约合并，\texttt{src/solver/alns.py}实现
ALNS和整数规划修复，\texttt{src/solver/technical\_stops.py}实现技术经停与加油
标签搜索，\texttt{src/solver/evaluator.py}实现内存路线评价，
\texttt{src/validation/solution\_validator.py}从最终CSV重新计算约束和五项指标。
统一运行入口为\texttt{scripts/04\_solve\_q1\_alns.py}。

程序保存\texttt{q1-routes.csv}、\texttt{q1-assignments.csv}、
\texttt{metrics.json}、\texttt{validator.json}、\texttt{q1-convergence.csv}和
\texttt{operator\_stats.csv}。最终方案覆盖全部1600人，内存指标和独立Validator
指标完全一致，且Validator未报告任何约束错误。

\subsection{结果与理论下界对照}

表~\ref{tab:q1-final-results}比较单设施底解、广义节约解和两阶段ALNS结果。由于本文
严格优先最小化总飞机使用时间，ALNS可以在一级目标下降时接受人员在途时间略有增加的
方案；这与式~\eqref{eq:q1-objective}规定的词典序目标一致。

\begin{table}[htbp]
  \centering
  \caption{问题一不同求解阶段的结果}
  \label{tab:q1-final-results}
  \resizebox{\linewidth}{!}{%
  \begin{tabular}{lrrrrr}
    \toprule
    方法 & 飞机时间/min & 人员在途/min & 架次数 & 燃油/kg & 利用率/\%\\
    \midrule
    单设施容量动态规划（B0） & 17222 & 120609 & 110 & 142601.0 & 43.047\\
    广义节约（B1）           & 15743 & 120396 & 100 & 128562.3 & 46.757\\
    两阶段ALNS（B2）         & \textbf{15418} & 121138 & \textbf{95}
                                  & \textbf{124811.4} & \textbf{48.696}\\
    \bottomrule
  \end{tabular}}
\end{table}

相对广义节约解，ALNS将总飞机使用时间减少325 min，改进率为
\begin{equation}
  I_{\mathrm{Q1}}
  =\frac{15743-15418}{15743}\times100\%
  =2.064\%.
\end{equation}
同时减少5个架次和3750.9 kg燃油，座位公里利用率提高1.939个百分点。最终95条路线
中，T1、T2和T3分别为4、46和45架次；平均每架运输16.842人，最低载客数由B1的8人
提高至12人，说明主要收益来自低利用率路线的重新分批。

采用放松容量共享、燃油和停靠次数约束得到的总飞机时间下界为13337 min，因而最终解
相对一级目标下界的间隙为
\begin{equation}
  \operatorname{Gap}_{\mathrm{Q1}}
  =\frac{15418-13337}{13337}\times100\%
  =15.603\%.
\end{equation}
该数值是相对全局最优解误差的保守上界，而不是算法实际误差的精确值。人员在途时间、
架次数和燃油消耗量的独立理论下界分别为95328 min、85架次和113878.73 kg；这些分量
下界是在不同放松条件下分别取得的，不能认为能够由同一方案同时达到。

\begin{figure}[htbp]
  \centering
  \includegraphics[width=0.86\linewidth]{figures/q1_alns_convergence.pdf}
  \caption{问题一两阶段ALNS的历史最优值收敛过程}
  \label{fig:q1-alns-final-convergence}
\end{figure}

图~\ref{fig:q1-alns-final-convergence}中虚线为总飞机使用时间理论下界，竖直点线为两阶段
搜索的分界。历史最优总飞机时间由15743 min逐步降至15418 min，架次数由100降至95，
说明扩大破坏规模能够进一步突破单纯路线合并形成的局部最优。
